\section{Quantum Mechanics Overview \label{ch2:quantumMechanicsOverview}}
    \begin{table}[H]
        \makebox[\linewidth][c]{
            \begin{tabular}{|c|c|c|c|c|c|c|}
                \hline
                Molecule & $\omega_e\Units{cm^{-1}}$ & $\omega_e
                x_e\Units{cm^{-1}}$ &  $\omega_e y_e\Units{cm^{-1}}$ &
                $B_e\Units{cm^{-1}}$ & $\alpha_e\Units{cm^{-1}}$ & ZPE
                $\Units{cm^{-1}}$ \\ 
                \hline\hline
                $NO+$ & $2376.72$ & $16.255$ & $-0.01562$ & $1.997195$ &
                $0.018790$ & $1184.33$ \\ \hline
                $N2+$ & $2207.0115$ & $16.0616$ & $-0.04289$ & $1.93176$ &
                $0.0181$ & $1099.40$ \\ \hline 
                $O2+$ & $1905.892$ & $16.489$ & $0.02057$ & $1.689824$ &
                $0.019363$ & $948.91$ \\ \hline
                $NO$ & $1904.1346$ & $14.08836$ & $0.01005$ & $1.704885$ &
                $0.0175416$ & $948.647$ \\ \hline
                $N2$ & $2358.57$ & $14.324$ & $-0.00226$ & $1.99824$ &
                $0.017318$ & $1175.78$ \\ \hline
                $O2$ & $1580.161$ & $11.95127$ & $0.0458489$ & $1.44562$ &
                $0.0159305$ & $787.380$ \\ \hline 
            \end{tabular}
        }
            \caption{Spectroscopy constants \cite{Irikura2007}\label{tab:spectroscopyConstant}}
    \end{table}

    Some authors refer the constants in \cref{tab:spectroscopyConstant}
    as the \textit{Dunham} coefficients \cite{Dunham1932} $\Parenthesis{Y_{ij}
    = Y_{vJ}}$ \cite{Chandra2005}. However, in the literature these
    coefficients can be called the \textit{Duhham} potential
    coefficients which are not the same. Ogilvie \cite{Ogilvie1976} called the
    $Y_{vJ}$ \textit{energy-coefficients} and $a_m$ the \textit{Dunham} potential
    coefficients, these will be the convention that we will use for the rest of
    this work.
    \Cref{eq:dunhamCoefficients} provides a relationship between the
    energy coefficients and the spectroscopy constants found in
    \cref{tab:spectroscopyConstant}.
    \begin{equation}\label{eq:dunhamCoefficients}
        Y_{10} = \omega_e,\; Y_{20} = -\omega_ex_e,\; Y_{30} = \omega_ey_e, \;
        Y_{01} = B_e,\; Y_{11} = -\alpha_e
    \end{equation}

    The \textit{Dunham} potential coefficients are found using
    \cref{eq:dunhamA0} \cite{Ogilvie1976}, \cref{eq:dunhamA1,eq:dunhamA2}
    \cite{Herschbach1961} and \cref{eq:dunhamA3}.
    \cite{Noorizadeh2004}.
    \begin{subequations}\label{eq:dunhamCoefficients}
        \begin{align}
            a_0 &= \frac{\omega_e^2}{4B_e} \label{eq:dunhamA0} \\
            a_1 &= -\Parenthesis{1 + \frac{\alpha_e\omega_e}{6B_e^2}}\label{eq:dunhamA1} \\
            a_2 &= \frac{5}{4}a_1^2 -
            \frac{2}{3}\Parenthesis{\frac{\omega_ex_e}{B_e}}\label{eq:dunhamA2} \\
            a_m &= \Parenthesis{\frac{12}{a_1}}^{m-2} \Parenthesis{2^{m+1}-1 }\Parenthesis{ \frac{a_2}{7}}^{m-1}\prod_{k=0}^{m-3}\frac{1}{m+2 – k} \label{eq:dunhamA3}
        \end{align}
    \end{subequations}


    \begin{table}[H]
        \makebox[\linewidth][c]{
            \begin{tabular}{|c|c|c|c|c|}
                \hline
                Molecule & $a_0\times 10^{5}$ & $a_1$ & $a_2$ & $a_3$ \\ 
                \hline\hline
                $NO+$ & $7.071$ & $-2.866$ & $4.842$ & $-6.009$ \\ \hline
                $N2+$ & $6.304$ & $-2.784$ & $4.146$ & $-4.536$ \\ \hline 
                $O2+$ & $5.374$ & $-3.154$ & $5.929$ & $-8.189$ \\ \hline 
                $NO$ & $5.317$ & $-2.915$ & $5.114$ & $-6.592$ \\ \hline 
                $N2$ & $6.960$ & $-2.705$ & $4.367$ & $-5.179$ \\ \hline
                $O2$ & $4.318$ & $-3.008$ & $5.795$ & $-8.205$ \\ \hline
            \end{tabular}
        }
            \caption{Potential \textit{Dunham} coefficients, \cref{eq:dunhamCoefficients}}\label{tab:potentialDunhamCoefficients}
    \end{table}

    
    The frequencies of transitions between electronic levels lie in the
    visible and ultraviolet, the vibrational transition in the infrared, and
    the rotation transition in the far infrared \cite{Schwabl2007}.
    \cref{fig:energySpectrum} shows the energy spectrum, this work will stay in
    the low visible light, so the electron transitions can be ignored. 

    \begin{figure}[H]
            \makebox[\textwidth][c]{
                \includegraphics[width=1.0\textwidth]{chapter2/energySpectrum}
            }
            \caption{Energy spectrum}\label{fig:energySpectrum}
    \end{figure}

    The \textit{Hamiltonian} of the system can be approximated using the
    \textit{Born}-\textit{Oppenheimer} approximation, which separates the
    motion of the nuclei and the motion of the electrons based on the fact that
    nuclei are much heavier and thus move much more slowly than electrons
    \cite{McNab2017}. The vibrational energy and rotational energy are calculated using
    \cref{eq:vibrationalEnergy} and \cref{eq:rotationalEnergy}. Many textbooks
    \todo{Cite textbooks} do not add the second term to these equations. These
    terms come from perturbation theory and are called the \textit{anharmonic}
    terms \todo{Find Reference}.
    
    \begin{subequations}\label{eq:vibrationalRotationalEnergies}
        \begin{align}
            \varepsilon_{trans} &= \label{eq:tranlationalEnergy} \\
            \varepsilon_{vib} &=
            \Parenthesis{\omega_e\Parenthesis{v + \frac{1}{2}}
                -\omega_ex_e\Parenthesis{v + \frac{1}{2}}^2} \UnitsFrac{1}{cm} \label{eq:vibrationalEnergy} \\
            \varepsilon_{rot} &= \Parenthesis{B_e -
            \alpha_e\Parenthesis{v +
            \frac{1}{2}}}J\Parenthesis{J+1}\UnitsFrac{1}{cm}\label{eq:rotationalEnergy}
        \end{align}
    \end{subequations}

    \todo{Figure where to add this \cite{Gupta2016}}

    The Boltzmann distribution, which describes the particle distribution over different energy states at thermal equilibrium, is applied to calculate the energy levels as a function of the vibrational quantum number. This distribution is given in \cref{eq:boltzmanDistribution}.
 
    \begin{equation}\label{eq:boltzmanDistribution}
        P(\varepsilon) = \frac{g_i e^{-\beta \mathcal{H}}}{Z}
    \end{equation}
 
    It should be noted that in a system of harmonic oscillators, the
    vibrational energy relaxes through a series of Boltzmann distributions.
    However, this is not the case for anharmonic oscillators
    \cite{Anketell1974}. Therefore, for this approach, it is assumed that the system contains only harmonic oscillators and the anharmonic $\Parenthesis{\omega_ex_e}$ term in vibrational energy equation REFEQ and the anharmonic $\Parenthesis{\alpha_e}$ in the rotational energy equation are set to zero.
